%==============================================================================
%  10-ai-by-function.tex — AI Usage Across Organizational Functions
%  (merged: performance change + function playbook + agent workflows)
%==============================================================================
\strategyPart{10 \textbar\ AI by Function}{AI Usage and Performance Change by Function}

AI's measured impact varies sharply by the type of work. The research allows
a function-by-function view of both adoption and performance change.

\section{Measured performance change by work type}

\begin{center}
  \renewcommand{\arraystretch}{1.5}
  \rowcolors{2}{inSnow}{white}
  \fitwidth{\begin{tabular}{>{\raggedright\arraybackslash}p{3.0cm}>{\raggedright\arraybackslash}p{4.4cm}>{\raggedright\arraybackslash}p{5.0cm}}
    \rowcolor{inNavy}
    \hcell{Function} & \hcell{AI application} & \hcell{Measured performance change} \\
    \textbf{Software engineering} & Coding assistants & \textbf{55.8\% faster} task completion in controlled trials [GitHub Copilot study, 2022]. \\
    \textbf{Customer support} & Generative-AI copilots & \textbf{+14\%} issues resolved per hour; \textbf{+34\%} for novice agents [MIT/NBER call-center field experiment, 2023]. \\
    \textbf{Knowledge work} & Copilots and assistants & \textbf{64\%} of workers report digital debt; AI frees time for higher-value work [Microsoft Work Trend Index 2023]. \\
    \textbf{Consulting / analysis} & LLM assistance & \textbf{25--40\%} faster on idea-generation tasks; \textbf{+40\%} quality on creative tasks in controlled experiments [BCG 2023]. \\
    \textbf{Marketing / sales} & Content generation, lead scoring & \textbf{75\% of gen-AI value} across customer operations, marketing/sales, engineering and R\&D [McKinsey 2023]. \\
  \end{tabular}}
\end{center}

\section{AI use by organizational function}

\vspace{-0.3em}
\infobox{\textbf{Customer experience}}{%
  Chatbots, personalization, intent analysis, contact-center copilots,
  complaint classification, journey orchestration. KPI: first-contact
  resolution, average handling time, CSAT, churn [WEF; CRM data].}

\infobox{\textbf{Operations \& supply chain}}{%
  Demand forecasting, inventory optimization, predictive maintenance, route
  optimization, quality control, digital twins. KPI: forecast accuracy,
  inventory cost, downtime [ERP/IoT data; WEF].}

\infobox{\textbf{Finance}}{%
  Fraud detection, cash-flow prediction, invoice processing, automated
  reporting, scenario planning. KPI: fraud loss, close time, error rate.}

\infobox{\textbf{Human resources}}{%
  Recruitment support, skills mapping, workforce planning, learning
  recommendations, attrition analysis. KPI: time-to-hire, attrition, internal
  mobility [WEF Future of Jobs; OECD].}

\infobox{\textbf{Marketing \& sales}}{%
  Lead scoring, segmentation, campaign generation, dynamic pricing, account
  research. KPI: conversion rate, sales-cycle duration, CAC.}

\infobox{\textbf{R\&D / product}}{%
  Generative design, simulation, experiment prioritization, requirements
  generation, product analytics. KPI: cycle time, experiment throughput.}

\infobox{\textbf{Strategic planning}}{%
  Scenario analysis, competitive intelligence, risk forecasting,
  capital-allocation support. KPI: forecast error, decision latency.}

\infobox{\textbf{Product management}}{%
  Customer research, discovery, prioritization, roadmap planning,
  documentation. KPI: time-to-insight, feature-adoption rate.}

\section{AI as an autonomous workflow agent}

An agent interprets a goal, retrieves data, plans actions, uses approved
tools, validates results, requests approval where necessary, executes and
logs. Every agent needs defined \textbf{identity, role, permissions, tool
scope, data scope, spending limits, escalation paths and logging} — Deloitte
reports organizations increasingly customize agents, making identity,
authorization, monitoring and human accountability essential
[Deloitte 2026; Gartner 2024].

\section{Three worked agent workflows}

\textbf{Customer-service agent:} customer request $\rightarrow$ intent
classification $\rightarrow$ history retrieval $\rightarrow$ policy check
$\rightarrow$ response generation $\rightarrow$ risk and quality validation
$\rightarrow$ human approval for sensitive cases $\rightarrow$ response
$\rightarrow$ audit record.

\textbf{Business-intelligence agent:} manager question $\rightarrow$ intent
agent $\rightarrow$ data-discovery agent finds trusted data $\rightarrow$
query agent builds and validates the query $\rightarrow$ analysis agent
calculates results $\rightarrow$ visualization agent charts them
$\rightarrow$ explanation agent writes the answer $\rightarrow$ human review
for important decisions.

\textbf{Product-management agent:} opportunity submitted $\rightarrow$
customer-research agent analyzes feedback $\rightarrow$ market agent analyzes
competitors $\rightarrow$ analytics agent reviews usage $\rightarrow$
feasibility agent identifies constraints $\rightarrow$ prioritization agent
scores the opportunity $\rightarrow$ \textbf{the product manager makes the
decision}.

These patterns support the professional without replacing strategic judgment
[Anthropic; Microsoft agent framework].

\section{Interpretation}

The productivity statistics are the strongest evidence for the problem-first
philosophy: the gains are concentrated where the problem was precisely defined
(a coding task, a customer interaction) and the workflow was designed for
human--AI collaboration. Where the problem is vague, the measured gains
disappear — which is exactly the population captured by the 30\% proof-of-
concept abandonment prediction [Gartner 2023].

\takeaway{Measured gains concentrate where the problem is precisely defined and the workflow is redesigned; adoption without redesign yields usage, not value.}
